\chapter{Computer addiction}
\index{Computer addiction}

\medskip
\noindent\textit{\textbf{Under review.} This chapter is an AI-assisted educational draft; it carries no source citations and is not a substitute for professional medical advice.}

\section*{Definition and Overview}
Computer addiction, also called problematic internet use or internet gaming disorder, is a pattern of behaviour in which a person loses control over their use of computers, the internet, gaming, or social media, and continues using them despite real harm to their life. It is not the same as simply spending a lot of time online. The World Health Organization recognises "gaming disorder" as a distinct condition when impaired control, increasing priority given to gaming, and continued use despite negative consequences are present for at least a year and cause significant distress or disability. The term computer addiction is often used more loosely to describe heavy, hard-to-control use of screens of any kind.

\section*{Epidemiology}
Problematic gaming is most common in adolescent and young adult males, with international studies suggesting that roughly 1--5\% of gamers meet criteria for gaming disorder. Rates of broader problematic internet use vary widely depending on how it is defined, with estimates of around 1--10\% of the population. People with depression, anxiety, attention-deficit hyperactivity disorder (ADHD), or autism spectrum conditions are at higher risk, and the prevalence has grown as screens have become part of everyday life.

\section*{Aetiology and Pathophysiology}
Computer addiction develops through a combination of features of the activity, the person, and their circumstances:

\begin{itemize}
    \item \textbf{Reward and conditioning:} games and social media deliver unpredictable rewards that trigger the brain's reward system, creating a strong pull to keep going
    \item \textbf{Escape and coping:} for many, heavy use is a way of escaping stress, boredom, loneliness, or negative feelings
    \item \textbf{Co-existing conditions:} depression, anxiety, ADHD, and autism spectrum conditions make problematic use more likely
    \item \textbf{Personality factors:} impulsivity, low self-esteem, and a tendency to seek immediate rewards are risk factors
    \item \textbf{Social circumstances:} limited social support, family conflict, or social difficulty at school or work increase reliance on online worlds
    \item \textbf{Game design:} features such as levelling up, loot boxes, and social pressure within games are deliberately engaging and can encourage excessive play
\end{itemize}

\section*{Clinical Features}
\begin{itemize}
    \item Loss of control: repeated failed attempts to cut down or stop use
    \item Preoccupation: thinking about the activity even when not using it
    \item Withdrawal symptoms: irritability, restlessness, or low mood when not able to use screens
    \item Tolerance: needing more time online to get the same satisfaction
    \item Neglect: falling behind at school or work, and dropping hobbies and face-to-face relationships
    \item Continued use despite harm: carrying on even when it damages health, grades, jobs, or relationships
    \item Escaping negative moods by going online
    \item Sleep disturbance from late-night use, and physical effects such as eye strain, headache, poor posture, and a sedentary lifestyle
    \item Lying to family or others about the amount of time spent online
\end{itemize}

\section*{Differential Diagnosis}
\begin{itemize}
    \item Heavy but non-problematic use that does not cause harm or impaired control
    \item Depression or anxiety, which may be causing escape-driven use rather than addiction itself
    \item Attention-deficit hyperactivity disorder, which can present with difficulty shifting attention
    \item Autism spectrum conditions, where intense interest in a subject is a feature of the condition
    \item Obsessive-compulsive disorder, when online behaviour has a ritualistic quality
    \item Substance use or gambling problems, which may coexist and require separate treatment
\end{itemize}

\section*{When to Seek Medical Care}
Consult a GP or mental health professional if screen use is interfering with school, work, sleep, or relationships, if attempts to cut down keep failing, or if mood, appetite, or energy are significantly affected. A GP can arrange a mental health plan and referral to appropriate services.

\textbf{Contact local emergency services or go to an emergency department if there are thoughts of suicide or self-harm.} You can also contact a local crisis service. Sudden withdrawal can cause intense distress, and any signs of a mental-health crisis should be taken seriously.

\section*{Diagnosis and Investigations}
\begin{itemize}
    \item \textbf{Clinical interview:} a detailed history of screen use, attempts to control it, and its effects on sleep, school, work, and relationships
    \item \textbf{Screening tools:} questionnaires such as the Internet Addiction Test (IAT) can help measure problem severity
    \item \textbf{Assessment of co-morbidities:} screening for depression, anxiety, ADHD, and substance use
    \item \textbf{Functional assessment:} understanding what the person is escaping from and what online activity provides
    \item \textbf{Family interview:} gathering information from parents or partners, especially for children and teenagers
\end{itemize}

\section*{Management}
\begin{itemize}
    \item \textbf{Psychoeducation:} helping the person and family understand how problematic use develops and what realistic recovery looks like
    \item \textbf{Cognitive behaviour therapy (CBT):} the most studied approach, focused on identifying triggers, building coping skills, and changing habits around screens
    \item \textbf{Setting limits:} agreed screen-time budgets, phone-free times and places, and structured replacement activities
    \item \textbf{Motivational interviewing:} helping the person build their own reasons and motivation to change
    \item \textbf{Family involvement:} improving communication and reducing conflict around screen use
    \item \textbf{Treating underlying conditions:} addressing depression, anxiety, or ADHD that may be driving the behaviour
    \item \textbf{Education and vocational support:} help with returning to school or work as use reduces
\end{itemize}

\section*{Complications}
\begin{itemize}
    \item Academic failure, school refusal, or job loss
    \item Sleep deprivation and daytime fatigue
    \item Physical inactivity, weight gain, and associated health problems
    \item Eye strain, headaches, and musculoskeletal pain from poor posture
    \item Social isolation and conflict within the family
    \item Financial problems from in-game spending
    \item Worsening of underlying depression, anxiety, or other mental health conditions
\end{itemize}

\section*{Prognosis and Outcomes}
Many people with problematic computer use improve substantially with structured limits, therapy, and treatment of underlying conditions, particularly when the person is motivated and has family support. Outcomes tend to be better in younger people whose habits have not been entrenched for long. Some people continue to struggle and may relapse during stress or major life changes, and ongoing relapse-prevention work is often needed.

\section*{Prevention and Self-Care}
\begin{itemize}
    \item Set clear, agreed screen-time limits and stick to them
    \item Keep screens out of the bedroom and switch them off well before bedtime
    \item Schedule offline activities, exercise, and face-to-face time with family and friends
    \item Use built-in screen-time and app limits on devices
    \item Model healthy screen use for children
    \item Address stress, sleep, and mood problems early rather than using screens to escape them
    \item Seek help early if use is starting to interfere with daily life
\end{itemize}
